\documentclass[11pt]{article}
\usepackage{amsmath,amssymb,amsthm,booktabs,hyperref}
\newtheorem{theorem}{Theorem}
\newtheorem{lemma}[theorem]{Lemma}
\newtheorem{proposition}[theorem]{Proposition}
\newtheorem{corollary}[theorem]{Corollary}
\title{The empty and singleton middle cases for quintic graphs}
\author{human\_compression project}
\date{}
\begin{document}
\maketitle

\paragraph{Purpose and attribution.}
We consider connected simple graphs on $32$ vertices, of degree $5$,
without triangles, with a fixed geodesic of length eight. The middle
defined below is always nonempty. When it has one vertex, the previously
reviewed isolation theorem~\cite{singleton} and a common row inequality
reduce the graph to sides of orders six and seven. The side of order six
has six possible incidence forms. An incidence bound then forces
eccentricity four at the middle of the geodesic, so every feasible
construction has global diameter eight.

The original girth-four, diameter-eight census belongs to Cambie,
Goedgebeur, Jooken and Van den Eede~\cite{cages}. The neighborhood
reduction and the singleton isolation theorem are retained proof legs.
The present assembly reuses their arguments concerning independent side
rows, degree sums and false twins. No priority or historical novelty is
claimed. This written argument imports no numerical census and performs
no graph, state, Source, TeX or Lean execution. It does not by itself
acquire a pending search family or pay a native record.

\section*{The common setup}

Let $p_0,p_1,\ldots,p_8$ be a fixed shortest path. Put
\[
 A=N(p_0),\qquad B=N(p_4),\qquad C=N(p_8),\qquad
 I=A\cup B\cup C,\qquad J=V(G)\setminus I,
\]
and let $F=G[J]$ and $K=J\setminus\{p_0,p_4,p_8\}$. Define
\[
 L=\{v\in K:N(v)\cap A\ne\varnothing\},\qquad
 R=\{v\in K:N(v)\cap C\ne\varnothing\},\qquad
 M=K\setminus(L\cup R).
\]
All these sets are defined in the complete graph. An absent edge in a
partial graph does not imply that the corresponding completed vertex
belongs to $M$.

We use the setup lemma in~\cite{singleton}. The three neighborhoods are
disjoint independent five-sets, with no edges between them. Consequently
$|I|=15$, $|J|=17$, $|K|=14$, and $F$ has five edges. The three marked
centers are isolated in $F$. The sets $L,R$ are disjoint, have size at
least five, and have no edge between them; a vertex of $B$ cannot meet
both sides. Each vertex of $A$ has exactly four side neighbors in $L$,
and each vertex of $C$ exactly four in $R$. Those neighbor sets are
independent. Every side vertex has a positive incidence with its own
endpoint neighborhood. If $D$ is the induced graph on a side of size
$n$, it follows that
\begin{equation}\label{degreebound}
 \Delta(D)\leq n-4.
\end{equation}
Indeed a row containing a specified vertex gives three further side
vertices that cannot neighbor it. These conclusions use the fixed
geodesic, without a global diameter upper bound. The lower bound five
for a side excludes a closed $K_{5,5}$ component only when the side
would have four vertices; no such component claim is made for size five.

\begin{lemma}[The retained obstruction on six vertices]\label{six}
Suppose that a side graph $D$ on six vertices has no triangles and has
maximum degree at most two. Suppose that five labeled independent
four-subsets, together with $b$ further labeled independent four-subsets,
supply its external incidences, and that every side vertex has total
degree five. The pairs
\[
 (e(D),b)=(1,2)\quad\hbox{and}\quad(e(D),b)=(3,1)
\]
are impossible.
\end{lemma}
\begin{proof}
This is the common side obstruction in the last two cases of the proof
of the reviewed singleton isolation theorem~\cite[Theorem~2]{singleton}.
Its incidence argument applies independently of that theorem's middle
vertex. For one edge, its endpoints need eight external incidences
between them, while each of the seven independent rows contains at most
one endpoint. For three edges, $D$ must be
$P_4\sqcup2K_1$, $P_3\sqcup P_2\sqcup K_1$, or $3P_2$.
The last graph has no independent four-set. In either of the first two,
every independent four-set contains an isolated vertex. That vertex
would lie in all six rows and have degree at least six. These are
exactly the retained obstructions.
\end{proof}

\begin{lemma}[An edge and the external rows]\label{rowbound}
Let $D$ be a side graph and let $\alpha(v),\beta(v)$ count incidences
with its five endpoint rows and $b$ port rows. Suppose all these rows
are independent subsets and every side vertex has degree five. For
each edge $uv$ of $D$,
\begin{equation}\label{rowedge}
 \beta(u)+\beta(v)\geq5-d_D(u)-d_D(v),\qquad
 d_D(u)+d_D(v)\geq5-b.
\end{equation}
If equality holds in the second inequality, every external row contains
exactly one endpoint of that edge.
\end{lemma}
\begin{proof}
No independent row contains both endpoints. Thus
$\alpha(u)+\alpha(v)\leq5$ and
$\beta(u)+\beta(v)\leq b$. Substituting
$\alpha(v)+\beta(v)=5-d_D(v)$ gives both inequalities. Equality in
the total bound says that all $5+b$ rows achieve their maximum one
incidence, proving the last assertion. This is the same independent-row
count used in Lemma~\ref{six}, expressed without separate side cases.
\end{proof}

\section*{The middle cannot be empty}

\begin{theorem}\label{nonempty}
Under the common hypotheses, $M\ne\varnothing$.
\end{theorem}
\begin{proof}
Assume that $M$ is empty. Then $K=L\sqcup R$, and all five edges of
$F$ lie within the sides. Every vertex of $B$ has $p_4$ and four side
neighbors. Its side neighbors lie in one side and form an independent
four-set. Write $B_L,B_R$ for the resulting port parts. Both are
nonempty: if one were empty, the corresponding side, endpoint
neighborhood and endpoint center would form a proper component.

For a side write $n$ for its size, $e$ for its number of internal edges,
and $b$ for the number of its $B$ ports. Summing degrees on that side
gives
\begin{equation}\label{emptybalance}
 5n=20+4b+2e,\qquad 1\leq b\leq4.
\end{equation}
The two side sizes sum to fourteen and are at least five. Equation
\eqref{emptybalance} makes each size even, so the only possibilities
are six and eight. This parity observation also excludes size five
directly.

On the side of size six, equation~\eqref{emptybalance} becomes
$e=5-2b$. Nonnegativity and $b\geq1$ leave $(e,b)=(3,1)$ or $(1,2)$.
The five endpoint rows and the $b$ port rows are precisely the
independent four-subsets of Lemma~\ref{six}; bound~\eqref{degreebound}
gives maximum internal degree two. Both alternatives are impossible.
This contradicts $M=\varnothing$.
\end{proof}

\section*{A singleton middle}

We now assume that $M=\{h\}$. The reviewed theorem
\cite[Theorem~2]{singleton} gives
\begin{equation}\label{isolation}
 d_F(h)=0,\qquad N(h)=B=N(p_4).
\end{equation}
We use this conclusion as a proof leg; the argument excluding all
positive values of $d_F(h)$ is not repeated. Put $X=\{p_4,h\}$.
These are nonadjacent false twins, and $F$ again has all five edges
inside $L$ and $R$.

Every $B$ port meets both vertices of $X$ and three vertices of one side.
Let $b_L,b_R$ be the port counts and $n_L,n_R,e_L,e_R$ the side sizes
and internal edge counts. Connectedness makes both port parts nonempty.
The degree equations are
\begin{equation}\label{singlebalance}
 n_L+n_R=13,\qquad b_L+b_R=5,\qquad e_L+e_R=5,
 \qquad 3b_D=5n_D-20-2e_D.
\end{equation}
The endpoint rows have four side neighbors and port rows have three;
both kinds are independent sets.

\begin{proposition}\label{sizes}
Up to reversing the marked geodesic, the initial necessary side data are
\[
\begin{array}{c|cc|cc|cc}
 &n_L&n_R&e_L&e_R&b_L&b_R\\\hline
 \text{I}&5&8&1&4&1&4\\
 \text{II}&6&7&2&3&2&3
\end{array}
\]
\end{proposition}
\begin{proof}
Each side has size at least five, so their unordered sizes are $(5,8)$
or $(6,7)$. On a side of size five, bound~\eqref{degreebound} makes its
internal graph a matching. It cannot contain two edges, since then its
independence number is at most three whereas an endpoint row has size
four. Thus $e\leq1$. Equation~\eqref{singlebalance} gives $3b=5-2e$:
$e=0$ is impossible by integrality, and $e=1$ gives $b=1$. The opposite
data are then $n=8,e=4,b=4$.

On a side of size six, equation~\eqref{singlebalance} gives
$3b=10-2e$. Since $0\leq e\leq5$, its integral possibilities have
$e=2$ or $e=5$. The latter gives $b=0$ and contradicts connectedness.
Thus $e=2,b=2$, and the opposite data are $n=7,e=3,b=3$.
When normalizing the smaller side to $L$, reverse every $p_i$ to
$p_{8-i}$ and interchange $A,C$ and $L,R$ together. For an ordered-root
application this reversal is part of the supplied identification.
\end{proof}

\paragraph{The internal side graphs.}
On five vertices with one edge there is just $K_2\sqcup3K_1$.
On six vertices with two edges there are
\[
 P_3\sqcup3K_1,\qquad 2K_2\sqcup2K_1.
\]
On seven vertices with three edges and no triangles there are
\[
 P_4\sqcup3K_1,\quad K_{1,3}\sqcup3K_1,\quad
 P_3\sqcup K_2\sqcup2K_1,\quad3K_2\sqcup K_1.
\]
On eight vertices with four edges and no triangles there are the following
nine possibilities:
\begin{equation}\label{fourshapes}
\begin{split}
 &P_5\sqcup3K_1,\quad K_{1,4}\sqcup3K_1,\quad T_5\sqcup3K_1,
 \quad C_4\sqcup4K_1,\\
 &P_4\sqcup K_2\sqcup2K_1,\quad K_{1,3}\sqcup K_2\sqcup2K_1,
 \quad2P_3\sqcup2K_1,\\
 &P_3\sqcup2K_2\sqcup K_1,\qquad4K_2.
\end{split}
\end{equation}
Here $T_5$ is the tree with degree sequence $(3,2,1,1,1)$.
For completeness, a connected graph with four edges and no triangles
is either one of the three trees on five vertices or $C_4$. If it is
disconnected, the edge partitions $3+1$, $2+2$, $2+1+1$ and $1+1+1+1$
give respectively the next two, one, one and one possibilities in
\eqref{fourshapes}. Isolates pad the order to eight. The smaller lists
follow by the same connected-component argument.

Thus there are nine potential pairs of side shapes in case I and
$2\cdot4=8$ in case II. The total seventeen counts potential side
patterns before the full incidence constraints. It is
neither a count of ordinary graphs nor an assertion that every pattern
has a realization. Different choices of rows or roots can describe the
same ordinary graph.

\begin{proposition}\label{feasibility}
Case I has no incidence realization. In case II the side of order six
must be $P_3\sqcup3K_1$. Thus only the four displayed side graphs of
order seven remain as potential partners.
\end{proposition}
\begin{proof}
In case I the size-five side has one edge and one port. Its edge has
endpoint degree sum two, whereas Lemma~\ref{rowbound} requires at least
$5-1=4$. Thus all nine initial patterns in that case are impossible.
On the size-six side of case II, the matching $2K_2\sqcup2K_1$ also
has edge degree sum two, while the required lower bound is $5-2=3$.
The other shape $P_3\sqcup3K_1$ survives this obstruction. These arguments
eliminate thirteen of the seventeen initial side patterns; they do not
count the incidence realizations of the remaining four.
\end{proof}

\section*{The side of order six}

\begin{proposition}\label{sixforms}
Write its path as $x-v-y$, with isolates $Q=\{q_1,q_2,q_3\}$, and
let $r$ be the number of endpoint rows containing $v$. Every feasible
incidence form has $r\in\{1,2,3\}$. Its two port rows are precisely
\[
 \begin{array}{ll}
 3-r\text{ rows}&\{v\}\cup Q',\quad Q'\subset Q,\ |Q'|=2,\\
 r-1\text{ rows}&\{x,y,q\},\quad q\in Q.
 \end{array}
\]
If $\beta(q)$ is the resulting port column count, its five endpoint
rows are precisely $r$ copies of $\{v\}\cup Q$ and, for each $q\in Q$,
$\beta(q)$ copies of $\{x,y\}\cup(Q\setminus\{q\})$.
Conversely every choice of these port rows gives a feasible side form
with those endpoint rows.
\end{proposition}
\begin{proof}
The endpoint degree sums of both path edges are three. Since there are
two port rows, equality holds in Lemma~\ref{rowbound}. Every one of the
seven external rows therefore meets each path edge exactly once. Its
intersection with the path is either $\{v\}$ or $\{x,y\}$.
The center $v$ has internal degree two, so its total external incidence
is three. Consequently, if $r$ endpoint rows contain it, exactly $3-r$
port rows contain it. There are only two port rows and its endpoint
incidence is positive, giving $1\leq r\leq3$ and the displayed types.

An endpoint row containing $v$ uses all three isolates; one containing
$x,y$ omits exactly one isolate. At an isolate $q$, degree five makes
the number of those omissions equal to $\beta(q)$. This gives the
asserted endpoint multiset. The port rows use a total of
$2(3-r)+(r-1)=5-r$ isolate incidences, so there are exactly five endpoint
rows. Conversely the degrees at $x,y$ are
$1+(5-r)+(r-1)=5$, the degree at $v$ is $2+r+(3-r)=5$, and an isolate
has degree $(5-\beta(q))+\beta(q)=5$. Since $\beta(q)\leq2$, each
isolate has positive endpoint incidence. All chosen rows are independent,
and the same is true of their displayed endpoint counterparts.
\end{proof}

Up to permuting isolates, ports, endpoint rows and the two path ends,
these are exactly six local forms:
\[
\begin{array}{c|l|l}
 r&\text{choices among isolates in the two port rows}
   &\text{sorted values of }\beta\text{ on }Q\\\hline
 1&\text{equal pairs}&(2,2,0)\\
 1&\text{distinct pairs}&(2,1,1)\\
 2&\text{singleton inside the pair}&(2,1,0)\\
 2&\text{singleton outside the pair}&(1,1,1)\\
 3&\text{equal singletons}&(2,0,0)\\
 3&\text{distinct singletons}&(1,1,0)
\end{array}
\]
For $r=1$ or $3$, equality versus inequality of the two subsets exhausts
their orbits on a three-set. For $r=2$, membership of the singleton in
the pair does so. The integer $r$ itself is invariant under these local
role-preserving isomorphisms because $v$ is the unique vertex of internal
degree two. This is a count of local incidence forms, not ordinary graphs.
In every form $x,y$ are false twins in the complete construction.

\section*{The exact incidence construction}

Choose case II, put $D_L=P_3\sqcup3K_1$ and choose $D_R$ from the four
displayed side graphs of order seven. Take disjoint sets $L,R$ of their
displayed sizes and disjoint
sets $A,C$ of five row vertices each. Partition a further five-set $B$
into parts $B_L,B_R$ of the displayed sizes. Finally add four vertices
$p_0,p_8,p_4,h$.

For each $a\in A$ choose an independent four-subset $S_a$ of $D_L$;
for each $c\in C$ choose an independent four-subset $S_c$ of $D_R$.
For each $b\in B_D$ choose an independent three-subset $T_b$ of $D_D$.
Rows are distinct vertices even when their chosen subsets coincide.
On the left, Proposition~\ref{sixforms} gives exactly the same choices
through its six local forms, so its row classification need not be
repeated in a later enumeration.
Write
\[
 \alpha_D(v)=|\{a\text{ in its endpoint rows}:v\in S_a\}|,
 \qquad
 \beta_D(v)=|\{b\in B_D:v\in T_b\}|.
\]
The complete column condition is
\begin{equation}\label{column}
 \alpha_D(v)=5-d_{D_D}(v)-\beta_D(v)\geq1
 \qquad(v\in D_D).
\end{equation}
In particular, it is not enough to match just the total number of
incidences. Every individual side vertex must have a positive incidence
with its endpoint neighborhood.

Join $p_0$ to every vertex of $A$, $p_8$ to every vertex of $C$, and each
of $p_4,h$ to all five vertices of $B$. Join endpoint rows and ports to
their chosen subsets, retain the edges of $D_L,D_R$, and add no other
edges. This specifies every adjacency.

\begin{theorem}[Necessity and converse]\label{normalform}
Every graph under the common hypotheses with $|M|=1$ has this incidence
form, up to the explicitly recorded endpoint reversal. Conversely every
choice satisfying~\eqref{column} constructs a connected simple graph on
$32$ vertices, of degree $5$, girth exactly $4$ and diameter exactly eight,
with an eight-geodesic
from $p_0$ through $p_4$ to $p_8$. Its middle is exactly $\{h\}$.
The global diameter assertion follows from the central distance argument
in Theorem~\ref{diameter} below.
\end{theorem}
\begin{proof}
For necessity, isolation~\eqref{isolation} and
Propositions~\ref{sizes}--\ref{feasibility} supply the complete central
and side data. The side graphs occur in the displayed lists because
they have the specified orders and edge counts and have no triangles. Every endpoint
row has four side neighbors and each port three; the neighbor subsets
are independent because the complete graph has no triangles. Summing
the actual neighbors of an individual side vertex gives
\eqref{column}, with positivity by the definition of $L,R$. Thus every
actual adjacency is among those prescribed.

For the converse, the vertex count is $4+13+15=32$. Simplicity follows
from using subsets and distinct row vertices. The four added vertices have
degree five. Each endpoint row has its endpoint and four side neighbors;
each port has the two core neighbors and three side neighbors. At a side
vertex, equation~\eqref{column} gives total degree five.

There is no triangle within $D_L$ or $D_R$. A triangle involving a row
vertex would require adjacent vertices inside one of its chosen
independent subsets. No triangle can contain two row vertices, since
$A\cup B\cup C$ is independent. The specified center neighborhoods
are also independent, and there are no other edges. The graph therefore
has no triangles.

Every side vertex meets an endpoint row. All endpoint rows on one side
are joined through their endpoint center, so that entire region is
connected. Each port has three side neighbors; both port parts are
nonempty and all ports meet the two core vertices. Thus both regions
and the core lie in one component.

There is a four-cycle. Indeed the five endpoint rows on either side
contain a total of $5\binom42=30$ occurrences of unordered pairs of side
vertices. There are at most $\binom82=28$ possible pairs. Some pair
occurs in two distinct rows and gives a four-cycle. Together with
simplicity and the absence of triangles, this proves girth exactly four.

The two core vertices have neighborhood exactly $B$. Their distance
from either endpoint is four: a path of that length uses an endpoint
row, a side vertex and a port. A shorter path cannot occur since
$N(p_0)=A$, $N(p_8)=C$, $N(X)=B$, and there are no edges inside
$A\cup B\cup C$. Removing $X$ separates the two regions. Every path
between the endpoints must therefore use a core vertex, and has length
at least eight. Left and right paths of length four through $p_4$ give
a path of length eight, so the marked distance is exactly eight and
$p_4$ is its midpoint.

The sets $A,B,C$ are exactly the specified center neighborhoods. The
vertices $p_0,p_4,p_8$ lie outside their union, as do $L,R,h$.
Condition~\eqref{column} makes every side vertex meet its own endpoint
neighborhood, while $h$ meets neither. Thus the recovered middle is
precisely $\{h\}$, as asserted.
\end{proof}

\section*{Central distances and the diameter}

The following elementary separator argument is useful beyond the
singleton case. If two nonadjacent vertices have equal open
neighborhoods, interchanging them is an automorphism. In particular
their distances to every other vertex are equal.

\begin{lemma}[The core separator]\label{separator}
Suppose $X$ is a nonempty set of nonadjacent false twins in a connected
graph, and $G-X$ has two regions with no edges between them. Suppose
their respective endpoints have distance four from every vertex of $X$.
For a vertex $v$ in one region and the endpoint $p$ of the other,
\[
 d(v,p)=d(v,x)+4\qquad\text{for every }x\in X.
\]
\end{lemma}
\begin{proof}
All distances $d(v,x)$ are equal because the core vertices are false
twins. Any $v$--$p$ path first meets a vertex of $X$ after at least that
distance, and its remaining part has length at least four. Conversely
a shortest path from $v$ to any fixed $x\in X$, followed by a path of
length four to $p$, gives the reverse inequality.
\end{proof}

\begin{theorem}\label{diameter}
For the singleton construction, the global diameter is eight if and
only if
\begin{equation}\label{eccentricity}
 \operatorname{ecc}(p_4)=\max_{v\in V(G)}d(p_4,v)\leq4.
\end{equation}
Every feasible construction satisfies this bound. Its graph center is
exactly $X=\{p_4,h\}$.
\end{theorem}
\begin{proof}
The two regions are
\[
 \{p_0\}\cup A\cup L\cup B_L,
 \qquad \{p_8\}\cup C\cup R\cup B_R.
\]
Every vertex outside $X$ belongs to one of them. If the global diameter
is eight, Lemma~\ref{separator} gives $d(v,p_4)+4\leq8$ for every
$v\notin X$. The other core vertex has distance two from $p_4$, so
\eqref{eccentricity} follows. Conversely~\eqref{eccentricity} and the
triangle inequality give
\[
 d(u,v)\leq d(u,p_4)+d(p_4,v)\leq8
\]
for all vertices. The two marked endpoints already have distance eight,
so this is the exact global diameter.

We now prove the eccentricity bound directly. Write $\beta(v)$ for a
side vertex's port incidence. If $uv$ is an internal side edge, then
$d_D(u)+d_D(v)\leq e(D)+1$: its incident edges are among all $e(D)$
edges, with $uv$ counted twice. Here $e(D)$ is two or three. If both
port incidences were zero, \eqref{rowedge} would give
$0\geq5-d_D(u)-d_D(v)\geq1$, a contradiction. Thus the port support
is a vertex cover of each side graph.

A vertex with positive port incidence is within two of $p_4$. A
nonisolated vertex without port incidence has a port-supported neighbor,
so is within three. An isolated vertex without port incidence has all
five endpoint rows as neighbors and is a false twin of its endpoint;
it is at distance four. This treats every side vertex.

The union $T$ of port supports on a side has size at least three. Every
endpoint row on the side of size six, being a four-set, meets $T$ and
is within three of $p_4$. On the side of size seven the same is true
unless the row avoids $T$. In that event $|T|=3$ and the row is its
four-vertex complement $S$. Every port has precisely $T$ as its
three-subset, so $T$ is independent; $S$ is independent as the endpoint
row. The three internal edges must all run between $T$ and $S$.
Consequently some vertex of $S$ neighbors $T$ and is within three of
$p_4$. The row itself is within four. Every endpoint row is therefore
within four in all cases. Ports, the other core vertex and the endpoints
already have distances one, two and four. This proves
\eqref{eccentricity} for every feasible construction.

Both core vertices have eccentricity four; their equality follows from
the twin automorphism. The diameter forces radius at least four.
Every $v\notin X$ has positive distance from $X$ and hence distance
at least five from the opposite endpoint, by the separator lemma.
No such $v$ is a center. The center is therefore exactly $X$.
\end{proof}

The upper bound therefore requires no distance filter in an exhaustive
enumeration of these feasible incidence forms. A single distance search
from $p_4$ would suffice as an independent computational check. The proof
of the automatic bound uses the feasibility eliminations: both surviving
sides have at most three internal edges. It is not inferred merely from
the preliminary degree equations.

The same separator argument applies to any previously stated stratum
in which every middle vertex is isolated in $F$: then
$X=\{p_4\}\cup M$ has common neighborhood $B$. Whenever that stratum
has global diameter eight, its center is exactly $X$. In particular the
already specified cases of two and three isolated middle vertices have
center sizes three and four, respectively. Together with the singleton
case, these strata are disjoint as ordinary graph classes, since center
cardinality is an intrinsic invariant. This observation does not weaken
the hypotheses or alter the statements of the earlier notes.

\section*{The two halves and ordinary identification}

Deleting the intrinsic center $X$ leaves two connected graphs $H_{14}$
and $H_{16}$, of orders
\[
 1+5+6+2=14,\qquad1+5+7+3=16.
\]
Their ports are precisely the vertices of internal degree three; every
other vertex retains degree five. Thus the port sets are intrinsic even
when all root and row labels are forgotten. The two halves cannot be
exchanged by an ordinary isomorphism because their orders differ.

\begin{proposition}\label{leftunrooted}
The six local forms in Proposition~\ref{sixforms} give exactly six
unrooted half graphs of order fourteen.
\end{proposition}
\begin{proof}
They cover every such half because the local classification already
identifies all side, row and port labels while retaining a marked
endpoint. Forgetting that endpoint cannot introduce additional classes;
it could only merge some. We distinguish all six intrinsically.

For a half $H$ with port set $P$, give $v$ the level
$1+d_H(v,P)$. This equals its distance to $X$ after gluing. The port
set is intrinsic, so the numbers of vertices at each level are unrooted
isomorphism invariants. For the six forms these numbers are
\[
\begin{array}{c|l|rrrr}
 r&\text{port choices on }Q&1&2&3&4\\\hline
 1&\text{equal pairs}&2&3&7&2\\
 1&\text{distinct pairs}&2&4&7&1\\
 2&\text{singleton inside pair}&2&5&5&2\\
 2&\text{singleton outside pair}&2&6&5&1\\
 3&\text{equal singletons}&2&3&6&3\\
 3&\text{distinct singletons}&2&4&6&2
\end{array}
\]
Here level one consists of the two ports. The path vertices with
positive $\beta$, together with the port-supported isolates, form
level two. A path vertex with zero $\beta$ has level three. Every one
of the five endpoint rows has level three, since its four-subset meets
the support of a three-subset on six vertices. Finally the marked
endpoint and the zero-$\beta$ isolates have level four. Substituting
the six values of $r$ and the displayed $\beta$ partitions gives the
table. The six vectors are different, so no two forms merge as unrooted
halves.
\end{proof}

Let $\mathcal H_{16}$ be the set of unrooted isomorphism classes of
halves of order sixteen admitting at least one of the feasible right-side
incidence presentations above. All such presentations, endpoints and
row labels must be identified when defining this set. It is not the
class of all graphs with three vertices of degree three and thirteen
of degree five; the incidence presentation is an additional condition.
Let $h_{16}=|\mathcal H_{16}|$, which has not been determined here.

\begin{theorem}\label{product}
The ordinary isomorphism classes of target graphs admitting the
singleton-middle presentation are in bijection with
\[
 \{\text{the six unrooted halves of order fourteen}\}
 \times\mathcal H_{16}.
\]
In particular their number is $6h_{16}$.
\end{theorem}
\begin{proof}
Given two such halves, add two nonadjacent core vertices and join each
to every port in both halves. There are two left ports and three right
ports, so the core vertices have degree five and each port regains
its two missing neighbors. The chosen incidence presentations give the
construction in Theorem~\ref{normalform}, whose diameter is eight by
Theorem~\ref{diameter}. This proves surjectivity onto the stated target
stratum and shows that every pair of halves is permitted.

Conversely an ordinary isomorphism preserves the center. Removing it
recovers the connected components of orders fourteen and sixteen and
hence their unrooted classes. This proves that two pairs producing the
same ordinary graph must have the same two classes. Isomorphisms of
the two halves preserve their degree-three port sets automatically.
They extend across any bijection of the two core vertices, since each
core vertex meets every port. There is no port alignment to choose
between the halves. Thus equal pairs also produce isomorphic graphs,
proving the asserted bijection and cardinality formula.
\end{proof}

The same argument gives
\[
 \operatorname{Aut}(G)\cong
 S_2\times\operatorname{Aut}(H_{14})\times\operatorname{Aut}(H_{16}).
\]
Restrictions to the intrinsic center and the distinguishable components
give an injective map, and every independent choice extends by the
complete core-to-port adjacencies. This statement is about ordinary
unrooted graph automorphisms. No numerical value of $h_{16}$, no sum
over the four right-side patterns, and no full ordinary census is
claimed here.

\begin{corollary}\label{nonbipartite}
A graph under the common hypotheses with $|M|=1$ is nonbipartite.
\end{corollary}
\begin{proof}
In a bipartite graph of degree five on $32$ vertices both color classes
have size sixteen. The three marked centers lie in one color class,
and $I$ occupies fifteen vertices of the other. The remaining vertex
$z$ of that color has no neighbor in $A\cup C$ and belongs to $M$.
All its five neighbors lie in $J$, so $d_F(z)=5$. A singleton middle
would give $z=h$, contradicting~\eqref{isolation}.
\end{proof}

\paragraph{Use in the unfinished census.}
Theorem~\ref{nonempty} gives an empty-family conclusion when the exact
geodesic setup and the universal condition $M=\varnothing$ are certified
in every selected completion. The singleton normal form supplies a
smaller finite incidence domain only when its exact hypotheses are
likewise acquired. A presently missing side edge is not a universal
exclusion. No pending family was inspected here, and no native return,
leaf count, reservation payment or ordinary class count is fabricated.
The seventeen initial side patterns reduce to four by the proved
incidence obstructions. The left incidences reduce to six unrooted
halves, and the global diameter needs no filter after feasibility.
The remaining right-side incidences and their unrooted identification
still require their own treatment before $h_{16}$ is known. Larger middle sets retain their own internal
incidence data; no isolation assertion for them is made here.

\begin{thebibliography}{9}
\bibitem{cages} Cambie, Goedgebeur, Jooken and Van den Eede,
\emph{On the order-diameter ratio of girth-diameter cages},
arXiv:2511.21144v3, Appendix~B, Table~6.
\bibitem{singleton} Human Compression project,
\emph{A singleton middle is necessarily isolated}, retained proof note,
Theorem~2 and its two final side obstructions.
The complete source is \texttt{single\_middle\_isolation.tex},
SHA256 \texttt{4381b9174a70afc5543aadfb63f9fe36394521d28558665685da42332e51a97a};
its independent mathematical review has SHA256
\texttt{f13282fd59f52f8296a0061be550350163ff2cc09ca07516f3af57f8abac49e7}.
\end{thebibliography}
\end{document}
